\chapter{Mycoplasma Genitalium PCR Test}

\section*{What Is This Test?}
Mycoplasma genitalium PCR detects genetic material from \textit{Mycoplasma genitalium}. This is a sexually transmitted bacterium. It can cause urethritis, cervicitis, pelvic inflammatory disease, and persistent or recurrent genitourinary symptoms. Molecular testing is preferred because the organism is difficult and slow to culture.

\section*{Alternative Names}
\begin{itemize}
  \item M. genitalium NAAT
  \item Mycoplasma Genitalium PCR
  \item Mgen Test
  \item Mycoplasma Genitalium Molecular Test
\end{itemize}
\section*{How It Is Performed}
The test uses a nucleic-acid amplification assay on first-void urine or a vaginal, cervical, urethral, or rectal swab. The choice depends on symptoms and exposure site. Some assays also detect macrolide-resistance mutations, which help guide treatment.

\section*{Why It Is Ordered}
\begin{itemize}
  \item Evaluating persistent or recurrent urethritis or cervicitis
  \item Assessing pelvic inflammatory disease when \textit{M. genitalium} is a possible cause
  \item Testing selected sexual contacts or symptomatic patients according to local guidance
  \item Resistance-guided management when a positive result is detected
\end{itemize}

\section*{Interpretation and Limitations}
A positive result means \textit{M. genitalium} nucleic acid was found. It must be read with symptoms because asymptomatic infection can occur. A negative result does not exclude infection at an untested anatomical site. Routine screening of asymptomatic people is not generally recommended in many guidelines. Treatment should follow current antimicrobial-resistance guidance. Testing for chlamydia, gonorrhoea, HIV, syphilis, and other infections may be indicated based on the clinical context.
\section*{Specimen and Preparation}
The laboratory may request first-void urine or a swab from the symptomatic or exposed anatomical site. Fasting is not required. Tell the clinician about recent antibiotics, pregnancy, persistent symptoms, and recent sexual contacts. These affect interpretation and treatment planning. A specimen from one site does not reliably assess another site.

\section*{Risks and Safety}
Urine collection has no procedural risk. A swab may cause brief discomfort or spotting. Do not self-treat with leftover antibiotics. Resistance is common, and inappropriate treatment can cause adverse effects and treatment failure. Partners and testing for other sexually transmitted infections should be addressed according to local guidance.

\section*{Understanding Results and Follow-up}
The report should state whether nucleic acid was detected and whether a macrolide-resistance marker was assessed. A positive result does not identify when the infection was acquired. It also does not prove that every symptom is caused by M. genitalium. Persistent symptoms may require adherence review, resistance testing, specialist advice, and a test of cure at the interval recommended by current guidance. Asymptomatic screening is generally not recommended.

\section*{Regulatory Status and Authorization Date}
This chapter describes a test category, not a specific commercial product. FDA, CDSCO, EU IVDR, and MHRA approval or registration dates therefore cannot be assigned without the assay manufacturer, product name, instrument, and intended use. For laboratory-developed tests, an individual agency approval date may not exist. See the Preface for the interpretation of regulatory status.

\clearpage
